%! Inverse Trigonometric Functions — Unit 3, Lesson 3.9 — Guided Notes
\documentclass[10pt]{article}
\usepackage{precalculus-article}
\usepackage{precalculus-boxes}
\newcommand{\TallMath}[1]{$\displaystyle #1\rule[-1.4em]{0pt}{3.2em}$}

\begin{document}

\pageheader{Unit 3, Lesson 3.9}{Guided Notes}
\namedateperiod

\vspace{0.10in}

%----------------------------------------------------------------------
\begin{notesbox}{LO 3.9.A \; Inverse Trigonometric Functions}

%--- Vocabulary row
\begin{skillbox}[Key Vocabulary]{greenbox}
\begin{multicols}{2}
\begin{itemize}[leftmargin=1.4em, itemsep=1pt]
    \item \termblanklong{arcsine}
    \item \termblanklong{arccosine}
    \item \termblanklong{arctangent}
    \item \termblanklong{restricted domain}
    \item \termblanklong{principal value}
\end{itemize}
\end{multicols}
\end{skillbox}

\vspace{0.08in}

%--- Section 1
\textbf{Section 1 --- Why Do We Need Restricted Domains?}

The table below shows selected values of $f(x) = \sin x$.

\begin{center}
\renewcommand{\arraystretch}{1.8}
\begin{tabular}{|c|c|c|c|c|c|c|c|}
\hline
$x$ &
$-\dfrac{\pi}{2}$ & $-\dfrac{\pi}{6}$ & $0$ &
$\dfrac{\pi}{6}$ & $\dfrac{\pi}{2}$ & $\dfrac{5\pi}{6}$ & $\pi$ \\
\hline
$\sin x$ & $-1$ & $-\tfrac{1}{2}$ & $0$ &
$\tfrac{1}{2}$ & $1$ & $\tfrac{1}{2}$ & $0$ \\
\hline
\end{tabular}
\end{center}

Notice that $\sin\!\bigl(\tfrac{\pi}{6}\bigr) = \sin\!\bigl(\tfrac{5\pi}{6}\bigr) = \tfrac{1}{2}$.
This shows $\sin x$ is \blank{3cm} (circle: one-to-one \;/\; NOT one-to-one).

\medskip
Because the sine function is periodic, it \blank{5.5cm} the horizontal line
test, so it \blank{4cm} have an inverse over its full domain.

\medskip
\textbf{Solution:} Restrict the domain so the function becomes one-to-one.

\vspace{0.08in}

%--- Section 2
\textbf{Section 2 --- Restricted Domains and the Inverse Trig Functions}

\begin{center}
\renewcommand{\arraystretch}{2.2}
\begin{tabular}{|l|c|c|c|c|}
\hline
\textbf{Function} & \textbf{Restricted Domain} & \textbf{Range} &
\textbf{Inverse Name} & \textbf{Notation} \\
\hline
Sine   & \blank{3cm} & $[-1,\,1]$  & Arcsine   & $\arcsin x$ or $\sin^{-1}x$ \\
\hline
Cosine & \blank{3cm} & $[-1,\,1]$  & Arccosine & $\arccos x$ or $\cos^{-1}x$ \\
\hline
Tangent & \blank{3cm} & $(-\infty,\infty)$ & Arctangent & $\arctan x$ or $\tan^{-1}x$ \\
\hline
\end{tabular}
\end{center}

\medskip
The \textbf{output} of an inverse trig function is an \blank{3cm} measure
(the \textbf{principal value}).

\vspace{0.08in}

%--- Section 3
\textbf{Section 3 --- Evaluating Inverse Trig Functions}

\textit{Ask:} ``Which angle \emph{in the restricted range} has the given trig value?''

\begin{center}
\renewcommand{\arraystretch}{2.4}
\begin{tabular}{|l|l|l|}
\hline
\textbf{Expression} & \textbf{Think:} & \textbf{Answer} \\
\hline
$\arcsin\!\left(\dfrac{1}{2}\right)$ &
  Which $\theta\in[-\tfrac{\pi}{2},\tfrac{\pi}{2}]$ has $\sin\theta=\tfrac{1}{2}$? &
  \blank{2.5cm} \\
\hline
$\arccos(0)$ &
  Which $\theta\in[0,\pi]$ has $\cos\theta=0$? &
  \blank{2.5cm} \\
\hline
$\arctan(1)$ &
  Which $\theta\in(-\tfrac{\pi}{2},\tfrac{\pi}{2})$ has $\tan\theta=1$? &
  \blank{2.5cm} \\
\hline
$\arcsin\!\left(-\dfrac{\sqrt{2}}{2}\right)$ &
  Which $\theta\in[-\tfrac{\pi}{2},\tfrac{\pi}{2}]$ has $\sin\theta=-\dfrac{\sqrt{2}}{2}$? &
  \blank{2.5cm} \\
\hline
$\arccos\!\left(-1\right)$ &
  Which $\theta\in[0,\pi]$ has $\cos\theta=-1$? &
  \blank{2.5cm} \\
\hline
\end{tabular}
\end{center}

\end{notesbox}

\vspace{0.08in}

\begin{practicebox}{Practice: Evaluate}
Evaluate each expression.  Write the exact answer in radians.

\renewcommand{\arraystretch}{1.8}
\begin{tabularx}{\linewidth}{@{} X X X @{}}
  1. $\arcsin(1) = $ \blank{2cm}
  &
  2. $\arccos\!\left(\dfrac{\sqrt{3}}{2}\right) = $ \blank{2cm}
  &
  3. $\arctan\!\left(-1\right) = $ \blank{2cm} \\[4pt]
  4. $\sin^{-1}\!\left(-\dfrac{1}{2}\right) = $ \blank{2cm}
  &
  5. $\cos^{-1}\!\left(-\dfrac{\sqrt{2}}{2}\right) = $ \blank{2cm}
  &
  6. $\tan^{-1}(0) = $ \blank{2cm}
\end{tabularx}
\end{practicebox}

\end{document}
